\documentclass[UTF8,a4paper,12pt]{ctexart}
\usepackage{geometry}
\usepackage{amsmath}
\usepackage{booktabs}
\usepackage{graphicx}
\usepackage{hyperref}
\usepackage{listings}
\usepackage{xcolor}

\geometry{left=2.5cm,right=2.5cm,top=2.5cm,bottom=2.5cm}
\hypersetup{colorlinks=true, linkcolor=blue, urlcolor=blue, citecolor=blue}

\lstset{
  basicstyle=\ttfamily\small,
  breaklines=true,
  frame=single,
  columns=fullflexible
}

\title{基于卷积神经网络的文本情感分类}
\author{姓名：\underline{\hspace{4cm}} \quad 学号：\underline{\hspace{4cm}}}
\date{\today}

\begin{document}
\maketitle

\begin{abstract}
本文完成 Rotten Tomatoes 电影评论句子的五分类情感分类任务。项目采用 TextCNN 作为主模型，通过不同窗口大小的一维卷积捕获句子中的局部 n-gram 情感特征，并使用最大池化获得句子级表示。为缓解小规模数据上随机词向量容易过拟合的问题，本文使用 GloVe 预训练词向量初始化 embedding，并尝试静态词向量通道与可训练词向量通道结合的 TextCNN。实验同时实现 FastText 作为 baseline，用于比较卷积结构对分类效果的影响。代码提供 Docker 环境，支持在 GPU 上完成训练、验证和测试，并自动保存实验指标与预测结果。
\end{abstract}

\section{任务描述}

本项目面向句子级文本情感分类。输入为已经分词的英文电影评论句子，输出为五个情感类别之一：
very negative、negative、neutral、positive、very positive。评价指标以准确率为主，同时报告 macro F1 作为类别不均衡情况下的补充指标。

\section{数据集}

实验使用课程提供的 Rotten Tomatoes 预处理数据。数据划分如表~\ref{tab:data} 所示。

\begin{table}[h]
\centering
\caption{数据集规模}
\label{tab:data}
\begin{tabular}{lr}
\toprule
划分 & 样本数 \\
\midrule
训练集 & 8544 \\
验证集 & 1101 \\
测试集 & 2210 \\
\bottomrule
\end{tabular}
\end{table}

句子平均长度约为 18 个 token，最长句子约 53 个 token。实验中统一将句子长度截断或补齐到 56。词表由训练前提供的 \texttt{tokens2id.csv} 读取，并额外加入 \texttt{<pad>} 与 \texttt{<unk>}。

\section{方法}

\subsection{TextCNN}

TextCNN 首先将输入 token 序列 $x_1, x_2, \dots, x_n$ 映射为词向量序列。实验中使用 GloVe 300d 预训练词向量进行初始化：
\[
E = [e_1, e_2, \dots, e_n], \quad e_i \in \mathbb{R}^d
\]

随后使用多个卷积核大小 $h \in \{2,3,4,5\}$ 的一维卷积提取局部特征：
\[
c_i = f(W \cdot E_{i:i+h-1} + b)
\]
其中 $f$ 为 ReLU 激活函数。对每个卷积通道进行最大池化：
\[
\hat{c} = \max_i c_i
\]
最后拼接不同卷积核得到的特征，经过 dropout 和全连接层输出五分类 logits。

\subsection{Baseline}

为了验证卷积结构的作用，本文实现 FastText 风格 baseline。该模型对句子中所有 token 的 embedding 取平均，然后直接输入线性分类器。该方法结构简单，可以作为非卷积模型的参考对比。

\section{实验设置}

主要超参数如表~\ref{tab:hyperparams} 所示。

\begin{table}[h]
\centering
\caption{主要实验超参数}
\label{tab:hyperparams}
\begin{tabular}{lr}
\toprule
参数 & 取值 \\
\midrule
embedding 维度 & 300 \\
TextCNN 卷积窗口 & 2, 3, 4, 5 \\
每种窗口卷积核数量 & 128 \\
dropout & 0.6 \\
batch size & 64 \\
学习率 & 0.0005 \\
最大训练轮数 & 40 \\
优化器 & AdamW \\
\bottomrule
\end{tabular}
\end{table}

训练过程中使用 early stopping、learning rate scheduler、gradient clipping 与 label smoothing。为了比较不同正则策略，实验保留了 accuracy 优先、较小单通道模型、类别权重模型三个 TextCNN 配置。

\begin{table}[h]
\centering
\caption{TextCNN 配置变体}
\label{tab:variants}
\begin{tabular}{ll}
\toprule
配置 & 说明 \\
\midrule
textcnn\_glove\_acc & GloVe + 双通道 TextCNN，accuracy 优先 \\
textcnn\_glove\_small & GloVe + 单通道 TextCNN，较小模型 \\
textcnn\_glove\_weighted & GloVe + 双通道 TextCNN + 类别权重，macro F1 优先 \\
\bottomrule
\end{tabular}
\end{table}

所有实验均在 Docker 容器中运行，服务器使用 NVIDIA GPU。运行命令示例如下：

\begin{lstlisting}[language=bash]
docker run --rm -it --gpus '"device=6"' \
  -v "$PWD/outputs:/workspace/outputs" \
  -v "$PWD/checkpoints:/workspace/checkpoints" \
  -v "$PWD/embeddings:/workspace/embeddings" \
  nlp-textcnn-sentiment bash scripts/run_textcnn_sweep.sh
\end{lstlisting}

\section{实验结果}

请在服务器运行实验后，将 \texttt{outputs/*metrics.json} 中的结果填入表~\ref{tab:results}。

\begin{table}[h]
\centering
\caption{测试集实验结果}
\label{tab:results}
\begin{tabular}{lrr}
\toprule
模型 & Accuracy & Macro F1 \\
\midrule
FastText baseline & 待填写 & 待填写 \\
TextCNN & 待填写 & 待填写 \\
\bottomrule
\end{tabular}
\end{table}

\section{结果分析}

从结果上看，FastText baseline 只使用词向量平均表示句子，能够捕获整体词汇情感倾向，但对局部短语结构建模能力较弱。TextCNN 通过不同窗口大小的卷积核提取短语级特征，例如否定结构、情感形容词短语等，因此通常能够获得更好的分类效果。

如果混淆矩阵显示 negative 与 very negative、positive 与 very positive 之间混淆较多，这是合理现象，因为相邻情感强度类别语义边界较模糊。相比之下，neutral 与强情感类别之间的错误更能反映模型是否真正理解情感线索。

\section{总结}

本文实现了一个完整的文本情感分类实验流程，包括数据加载、TextCNN 建模、baseline 对比、训练验证、测试评估、Docker 部署和结果保存。实验代码结构清晰，可在服务器容器环境中复现。后续可以进一步引入预训练词向量、BERT 类预训练模型或更细致的超参数搜索，以继续提升分类效果。

\end{document}
